% Generated by exam_parser.latex_renderer from gaokao_2024_national_A_wen__1_.tex (schema 1.0).
% Compile with XeLaTeX or LuaLaTeX:  xelatex <this file>
\documentclass[UTF8,12pt]{ctexart}

% ---------- Packages ----------
\usepackage{amsmath,amssymb,bm}
\usepackage[a4paper,top=2.8cm,bottom=2.5cm,left=2.4cm,right=2.4cm]{geometry}
\usepackage{xcolor}
\usepackage{enumitem}      % questions / sub-questions / statements
\usepackage{titlesec}      % section headings
\usepackage{fancyhdr}      % running header and footer
\usepackage{etoolbox}      % keep stems together with their options
\usepackage{tasks}         % multiple-choice options
\usepackage{booktabs}      % table rules
\usepackage{tcolorbox}     % boxed reference material
\usepackage{tikz}          % figure placeholders

% ---------- Fonts ----------
% xeCJK does not treat circled digits (① ② ...) as CJK by default, so they
% would come from the Latin font and be missing. Use the CJK font instead.
\ifdefined\xeCJKsetcharclass\xeCJKsetcharclass{"2460}{"24FF}{1}\fi

% ---------- Exam metadata ----------
\newcommand{\examtitle}{2024 年普通高等学校招生全国统一考试}
\newcommand{\examsubtitle}{数学（文科） · 全国甲卷}
\newcommand{\examheaderleft}{2024 全国甲卷（文科数学）}
\newcommand{\examheaderright}{普通高等学校招生全国统一考试}

% ---------- Colours ----------
\definecolor{accent}{HTML}{1F4E79}
\definecolor{softgray}{HTML}{F2F4F7}
\definecolor{midgray}{HTML}{8A94A6}

% ---------- Page style ----------
\pagestyle{fancy}
\fancyhf{}
\setlength{\headheight}{24.1pt}
\renewcommand{\headrulewidth}{0.4pt}
\renewcommand{\headrule}{\hbox to\headwidth{\color{accent}\leaders\hrule height \headrulewidth\hfill}}
\fancyhead[L]{\small\color{midgray}\examheaderleft}
\fancyhead[R]{\small\color{midgray}\examheaderright}
\fancyfoot[C]{\small\color{midgray}-- \thepage\ --}

\setlength{\parindent}{0pt}
\setlength{\parskip}{0.5em}
\linespread{1.25}
\renewcommand{\arraystretch}{1.3}   % roomier tables

% ---------- Section headings ----------
\titleformat{\section}[block]
  {\Large\bfseries\color{accent}}{}{0pt}{}
  [{\color{accent}\titlerule[1.2pt]}]
\titlespacing*{\section}{0pt}{1.6em}{0.8em}

% ---------- Lists ----------
\newlist{questions}{enumerate}{1}
\setlist[questions]{label=\textbf{\color{accent}\arabic*.},
  leftmargin=2.2em, itemsep=1.1em, topsep=0pt, parsep=0.4em}

\newlist{subquestions}{enumerate}{2}
\setlist[subquestions,1]{label=(\arabic*), leftmargin=2.6em,
  itemsep=0.3em, topsep=0.2em, parsep=0pt}
\setlist[subquestions,2]{label=(\alph*), leftmargin=2.4em,
  itemsep=0.2em, topsep=0.1em, parsep=0pt}

\newlist{statements}{description}{1}   % ① ② ③ ④ propositions
\setlist[statements]{font=\normalfont, align=left, leftmargin=3em,
  labelwidth=1.6em, labelsep=0.4em, itemsep=0.1em, topsep=0.2em, parsep=0pt}

% never separate a stem from its options / sub-questions across pages
\BeforeBeginEnvironment{subquestions}{\nopagebreak}
\BeforeBeginEnvironment{tasks}{\nopagebreak}

% ---------- Multiple-choice options ----------
\settasks{label=(\Alph*), label-width=1.8em, column-sep=1em,
  before-skip=0.2em, after-skip=0.2em, after-item-skip=0.2em}

% ---------- Boxed reference material ----------
\newtcolorbox{referencebox}{colback=softgray, colframe=accent!40,
  boxrule=0.5pt, arc=2pt, left=1em, right=1em, top=0.4em, bottom=0.4em}

% ---------- Macros ----------
\newcommand{\blank}{\underline{\hspace{4em}}}                % fill-in blank
\newcommand{\choiceparen}{（\hspace{2em}）}                    % answer bracket
\newcommand{\boxlabel}[1]{{\bfseries\color{accent}#1}\quad}   % label inside a box
\newcommand{\examfigure}[3]{%   width (cm), height (cm), label
  \begin{tikzpicture}[baseline=(current bounding box.center)]
    \draw[dashed,midgray,thick,rounded corners=4pt,fill=softgray]
      (0,0) rectangle (#1,#2);
    \node[midgray,align=center,font=\small] at (#1/2,#2/2) {\textit{[ #3 ]}};
  \end{tikzpicture}}
\newcommand{\makeexamtitle}{%
  \begin{center}
    {\color{accent}\rule{\linewidth}{1.6pt}}\par\vspace{0.8em}
    {\LARGE\bfseries\examtitle}\par\vspace{0.4em}
    {\Large\bfseries\color{accent}\examsubtitle}\par\vspace{0.6em}
    {\color{accent}\rule{\linewidth}{0.5pt}}
  \end{center}}

\begin{document}

\makeexamtitle

%==================================================================
\section*{一、单选题}
%==================================================================

\begin{questions}

  % Question 1
  \item 若 $z = \sqrt{2}\,i$，则 $z\overline{z} =$ \choiceparen{}
    \begin{tasks}(4)
      \task $-2$
      \task $-\sqrt{2}$
      \task $\sqrt{2}$
      \task $2$
    \end{tasks}

  % Question 2
  \item 已知集合 $A = \{1, 2, 3, 4, 5, 9\}$，$B = \{x \mid x + 1 \in A\}$，则
    $A \cap B =$ \choiceparen{}
    \begin{tasks}(2)
      \task $\{1,2,3\}$
      \task $\{3,4,9\}$
      \task $\{1,2,3,4\}$
      \task $\{2,3,4,5\}$
    \end{tasks}

  % Question 3
  \item 设向量 $\boldsymbol{a} = (x + 1, x)$，$\boldsymbol{b} = (x, 2)$，%
    $\boldsymbol{c} = (1, 1)$，若 $\boldsymbol{a} - 2\boldsymbol{b}$ 与 $\boldsymbol{c}$
    共线，则 $x =$ \choiceparen{}
    \begin{tasks}(4)
      \task $\dfrac{5}{2}$
      \task $\dfrac{1}{2}$
      \task $-\dfrac{1}{2}$
      \task $-\dfrac{5}{2}$
    \end{tasks}

  % Question 4
  \item 某独唱比赛的决赛阶段共有甲、乙、丙、丁四人参加，每人出场一次，
    出场次序由随机抽签确定。则丙不是第一个出场，且甲或乙最后出场的概率是\choiceparen{}
    \begin{tasks}(4)
      \task $\dfrac{1}{6}$
      \task $\dfrac{1}{4}$
      \task $\dfrac{1}{3}$
      \task $\dfrac{1}{2}$
    \end{tasks}

  % Question 5
  \item 记 $S_{n}$ 为等差数列 $\{a_{n}\}$ 的前 $n$ 项和。已知 $S_{9}=1$，则
    $a_{3}+a_{7}=$ \choiceparen{}
    \begin{tasks}(4)
      \task $\dfrac{1}{9}$
      \task $\dfrac{2}{9}$
      \task $\dfrac{1}{3}$
      \task $\dfrac{2}{3}$
    \end{tasks}

  % Question 6
  \item 已知双曲线的两个焦点分别为 $(0,4)$，$(0,-4)$，点 $(-6,4)$ 在该双曲线上，
    则该双曲线的离心率为\choiceparen{}
    \begin{tasks}(4)
      \task $4$
      \task $3$
      \task $2$
      \task $\sqrt{2}$
    \end{tasks}

  % Question 7
  \item 设函数 $f(x)=\dfrac{e^{x}+2\sin x}{1+x^{2}}$，则曲线 $y=f(x)$ 在点 $(0,1)$
    处的切线与两坐标轴所围成的三角形的面积为\choiceparen{}
    \begin{tasks}(4)
      \task $\dfrac{1}{6}$
      \task $\dfrac{1}{3}$
      \task $\dfrac{1}{2}$
      \task $\dfrac{2}{3}$
    \end{tasks}

  % Question 8
  \item 函数 $y = -x^{2} + \left(e^{x} - e^{-x}\right)\sin x$ 在区间 $[-2.8, 2.8]$
    的图象大致为\choiceparen{}
    \begin{tasks}(2)
      \task \examfigure{4.6}{3.2}{图象 A}
      \task \examfigure{4.6}{3.2}{图象 B}
      \task \examfigure{4.6}{3.2}{图象 C}
      \task \examfigure{4.6}{3.2}{图象 D}
    \end{tasks}

  % Question 9
  \item 直线 $2x - y - 2 = 0$ 与圆 $x^{2} + y^{2} - 6x - 8y = 0$ 交于 $A, B$ 两点，则
    $|AB| =$ \choiceparen{}
    \begin{tasks}(4)
      \task $4$
      \task $5$
      \task $8$
      \task $10$
    \end{tasks}

  % Question 10
  \item 已知 $\dfrac{\cos\alpha}{\cos\alpha-\sin\alpha}=\sqrt{3}$，则
    $\tan\left(\alpha+\dfrac{\pi}{4}\right)=$ \choiceparen{}
    \begin{tasks}(4)
      \task $2\sqrt{3}+1$
      \task $2\sqrt{3}-1$
      \task $\dfrac{\sqrt{3}}{2}$
      \task $1-\sqrt{3}$
    \end{tasks}

  % Question 11
  \item 设 $\alpha, \beta$ 为两个平面，$m, n$ 为两条直线，且 $\alpha \cap \beta = m$。
    下述四个命题：

    \begin{statements}
      \item[①] 若 $m \parallel n$，则 $n \parallel \alpha$ 或 $n \parallel \beta$；
      \item[②] 若 $m \perp n$，则 $n \perp \alpha$ 或 $n \perp \beta$；
      \item[③] 若 $n \parallel \alpha$ 且 $n \parallel \beta$，则 $m \parallel n$；
      \item[④] 若 $n$ 与 $\alpha, \beta$ 所成的角相等，则 $m \perp n$。
    \end{statements}

    其中所有真命题的编号是\choiceparen{}

    \begin{tasks}(4)
      \task ①③
      \task ②④
      \task ①②③
      \task ①③④
    \end{tasks}

  % Question 12
  \item 记 $\triangle ABC$ 的内角 $A, B, C$ 的对边分别为 $a, b, c$，已知
    $B = 60^{\circ}$，$b^{2} = \dfrac{9}{4}ac$，则 $\sin A + \sin C =$ \choiceparen{}
    \begin{tasks}(4)
      \task $\dfrac{3}{2}$
      \task $\sqrt{2}$
      \task $\dfrac{\sqrt{7}}{2}$
      \task $\dfrac{\sqrt{3}}{2}$
    \end{tasks}

\end{questions}

%==================================================================
\section*{二、填空题}
%==================================================================

\begin{questions}[start=13]

  % Question 13
  \item 函数 $f(x) = \sin x - \sqrt{3} \cos x$ 在区间 $[0, \pi]$ 的最大值为 \blank。

  % Question 14
  \item 已知圆台甲、乙的上底面半径均为 $r_1$，下底面半径均为 $r_2$，圆台甲、
    乙的母线长分别为 $2(r_2 - r_1)$，$3(r_2 - r_1)$，则圆台甲与乙的体积之比为 \blank。

  % Question 15
  \item 已知 $a > 1$ 且 $\dfrac{1}{\log_8 a} - \dfrac{1}{\log_a 4} = -\dfrac{5}{2}$，则
    $a =$ \blank。

  % Question 16
  \item 当 $x > 0$ 时，曲线 $y = x^{3} - 3x$ 与曲线 $y = -(x - 1)^{2} + a$ 有两个交点，
    则 $a$ 的取值范围是 \blank。

\end{questions}

%==================================================================
\section*{三、解答题}
%==================================================================

\begin{questions}[start=17]

  % Question 17
  \item 记 $S_{n}$ 为等比数列 $\{a_{n}\}$ 的前 $n$ 项和，已知 $2S_{n}=3a_{n+1}-3$。
    \begin{subquestions}
      \item 求 $\{a_{n}\}$ 的通项公式；
      \item 求数列 $\{S_{n}\}$ 的前 $n$ 项和。
    \end{subquestions}

  % Question 18
  \item 某工厂进行生产线智能化升级改造。升级改造后，从该工厂甲、
    乙两个车间的产品中随机抽取 150 件进行检验，数据如下：

    \begin{center}
      \begin{tabular}{lcccc}
        \toprule
               & 优级品 & 合格品 & 不合格品 & 总计 \\
        \midrule
        甲车间 & 26     & 24     & 0        & 50 \\
        乙车间 & 70     & 28     & 2        & 100 \\
        \midrule
        总计   & 96     & 52     & 2        & 150 \\
        \bottomrule
      \end{tabular}
    \end{center}

    \begin{subquestions}
      \item 填写如下列联表：

        \begin{center}
          \begin{tabular}{lcc}
            \toprule
                   & 优级品       & 非优级品 \\
            \midrule
            甲车间 & \hspace{3em} & \hspace{3em} \\
            乙车间 &              &  \\
            \bottomrule
          \end{tabular}
        \end{center}

        能否有 $95\%$ 的把握认为甲、乙两车间产品的优级品率存在差异？能否有 $99\%$
          的把握认为甲、乙两车间产品的优级品率存在差异？
      \item 已知升级改造前该工厂产品的优级品率 $p = 0.5$。设 $\overline{p}$
        为升级改造后抽取的 $n$ 件产品的优级品率，如果
        $\overline{p} > p + 1.65\sqrt{\dfrac{p(1 - p)}{n}}$，
        则认为该工厂产品的优级品率提高了。根据抽取的 150 件产品的数据，
        能否认为生产线智能化升级改造后，该工厂产品的优级品率提高了？%
        $(\sqrt{150} \approx 12.247)$
    \end{subquestions}

    \begin{referencebox}
      \boxlabel{附：} $\displaystyle K^{2} = \frac{n(ad - bc)^{2}}{(a + b)(c + d)(a + c)(b + d)}$
      \begin{center}
        \begin{tabular}{cccc}
          \toprule
          $P(K^{2} \geqslant k)$ & 0.050 & 0.010 & 0.001 \\
          \midrule
          $k$                    & 3.841 & 6.635 & 10.828 \\
          \bottomrule
        \end{tabular}
      \end{center}
    \end{referencebox}

  % Question 19
  \item 如图，在以 $A, B, C, D, E, F$ 为顶点的五面体中，四边形 $ABCD$ 与四边形 $ADEF$
    均为等腰梯形，$EF \parallel AD$，$BC \parallel AD$，$AD = 4$，$AB = BC = EF = 2$，%
    $ED = \sqrt{10}$，$FB = 2\sqrt{3}$，$M$ 为 $AD$ 的中点。

    \begin{center}
      \examfigure{6}{4.2}{题 19 图：五面体 $ABCDEF$}
      % real figure: \includegraphics[width=6cm]{imgs/img_in_image_box_2059_1248_2291_1465.jpg}
    \end{center}

    \begin{subquestions}
      \item 证明：$BM \parallel$ 平面 $CDE$；
      \item 求 $M$ 到平面 $FAB$ 的距离。
    \end{subquestions}

  % Question 20
  \item 已知函数 $f(x) = a(x - 1) - \ln x + 1$。
    \begin{subquestions}
      \item 讨论 $f(x)$ 的单调性；
      \item 设 $a \leqslant 2$，证明：当 $x > 1$ 时，$f(x) < \mathrm{e}^{x-1}$。
    \end{subquestions}

  % Question 21
  \item 设椭圆 $C$：$\dfrac{x^{2}}{a^{2}}+\dfrac{y^{2}}{b^{2}}=1$ $(a > b > 0)$
    的右焦点为 $F$，点 $M\left(1,\dfrac{3}{2}\right)$ 在 $C$ 上，且 $MF \perp x$ 轴。
    \begin{subquestions}
      \item 求 $C$ 的方程；
      \item 过点 $P(4,0)$ 的直线交 $C$ 于 $A, B$ 两点，$N$ 为线段 $FP$ 的中点，直线 $NB$
        交直线 $MF$ 于点 $Q$，证明：$AQ \perp y$ 轴。
    \end{subquestions}

  % Question 22
  \item 在直角坐标系 $xOy$ 中，以坐标原点为极点，$x$ 轴正半轴为极轴建立极坐标系，曲线
    $C$ 的极坐标方程为 $\rho = \rho \cos \theta + 1$。
    \begin{subquestions}
      \item 写出 $C$ 的直角坐标方程；
      \item 设直线 $l$：$\begin{cases}x=t,\\ y=t+a\end{cases}$（$t$ 为参数），若 $C$ 与
        $l$ 相交于 $A, B$ 两点，且 $|AB|=2$，求 $a$。
    \end{subquestions}

  % Question 23
  \item 已知实数 $a, b$ 满足 $a + b \geqslant 3$。
    \begin{subquestions}
      \item 证明：$2a^{2}+2b^{2}>a+b$；
      \item 证明：$\left|a-2b^{2}\right|+\left|b-2a^{2}\right|\geqslant 6$。
    \end{subquestions}

\end{questions}

\end{document}
